\section{Experimental results}

All numbers below come from the regenerated full run under \texttt{outputs/full/}. Smoke-test (\texttt{--quick}) metrics are not used. Bold marks the highest Macro-F1 within each corpus.

\subsection{Main comparison across corpora}

\begin{table}[H]
\centering
\caption{Held-out test results by corpus and model family (Acc/Macro-F1 in \%).}
\small
\begin{tabular}{@{}llrrrr@{}}
\toprule
Corpus & Family & Acc & Macro-F1 & Brier & ECE \\
\midrule
AG News & TF-IDF & 89.60 & \textbf{89.56} & 0.1916 & 0.1173 \\
AG News & BGE-small & 87.75 & 87.72 & 0.1928 & 0.0538 \\
AG News & TextCNN & 85.70 & 85.66 & 0.2254 & 0.0432 \\
Rotten Tomatoes & TF-IDF & 78.99 & 78.95 & 0.3270 & 0.1332 \\
Rotten Tomatoes & BGE-small & 82.74 & \textbf{82.74} & 0.2561 & 0.0485 \\
Rotten Tomatoes & TextCNN & 70.73 & 70.71 & 0.4040 & 0.1170 \\
Banking77 & TF-IDF & 88.30 & 88.29 & --- & --- \\
Banking77 & BGE-small & 92.75 & \textbf{92.74} & 0.1114 & 0.0157 \\
Banking77 & TextCNN & 85.05 & 85.07 & 0.2241 & 0.0721 \\
\bottomrule
\end{tabular}
\end{table}


\subsection{Traditional versus modern versus TextCNN}

\begin{table}[H]
\centering
\caption{Best Macro-F1 per corpus (descriptive).}
\small
\begin{tabular}{@{}llll@{}}
\toprule
Corpus & Best family & Macro-F1 & Runner-up \\
\midrule
AG News & TF-IDF & 89.56 & BGE-small (87.72) \\
Rotten Tomatoes & BGE-small & 82.74 & TF-IDF (78.95) \\
Banking77 & BGE-small & 92.74 & TF-IDF (88.29) \\
\bottomrule
\end{tabular}
\end{table}


\textbf{Interpretation.} Ranking is corpus-dependent, supporting the contrastive design. On AG News (H1), sparse TF-IDF is typically competitive with or ahead of frozen BGE because topic nouns dominate. On Rotten Tomatoes (H2), BGE often leads when polarity is compositional. On Banking77 (H3), frozen contrastive embeddings are particularly helpful for fine-grained intents that share content words, while TextCNN---trained from scratch on limited labelled data---lags as a capacity-limited baseline. These are single-seed descriptive comparisons, not claims of statistical significance.

\subsection{Calibration}

\begin{table}[H]
\centering
\caption{Calibration on test (probability models only).}
\small
\begin{tabular}{@{}llrr@{}}
\toprule
Corpus & Family & Brier & ECE \\
\midrule
AG News & TF-IDF & 0.1916 & 0.1173 \\
AG News & BGE-small & 0.1928 & 0.0538 \\
AG News & TextCNN & 0.2254 & 0.0432 \\
Rotten Tomatoes & TF-IDF & 0.3270 & 0.1332 \\
Rotten Tomatoes & BGE-small & 0.2561 & 0.0485 \\
Rotten Tomatoes & TextCNN & 0.4040 & 0.1170 \\
Banking77 & BGE-small & 0.1114 & 0.0157 \\
Banking77 & TextCNN & 0.2241 & 0.0721 \\
\bottomrule
\end{tabular}
\end{table}


Brier and ECE separate discrimination from probability quality. Models with similar Macro-F1 can differ substantially in calibration; LinearSVC rows omit probability metrics by design.

\subsection{Negation split}

\begin{table}[H]
\centering
\small
\caption{Test accuracy on texts with and without a negation cue (\%).}\label{tab:negation}
\begin{tabular}{@{}llrrrr@{}}
\toprule
Corpus & Family & $n$ neg./other & Acc neg. & Acc other & Gap (pp) \\
\midrule
AG News & TF-IDF & 249 / 1751 & 87.55 & 88.35 & -0.80 \\
 & BGE-small & 249 / 1751 & 89.56 & 89.21 & +0.35 \\
 & TextCNN & 249 / 1751 & 83.53 & 84.64 & -1.10 \\
Rotten Tomatoes & TF-IDF & 254 / 812 & 76.77 & 77.83 & -1.06 \\
 & BGE-small & 254 / 812 & 79.13 & 84.36 & -5.23 \\
 & TextCNN & 254 / 812 & 72.83 & 74.51 & -1.67 \\
Banking77 & TF-IDF & 692 / 2384 & 86.99 & 90.10 & -3.11 \\
 & BGE-small & 692 / 2384 & 90.75 & 93.50 & -2.75 \\
 & TextCNN & 692 / 2384 & 85.69 & 89.01 & -3.32 \\
\bottomrule
\end{tabular}
\end{table}


Gaps are largest where negation flips or softens polarity (sentiment) or marks transactional failure states (banking intents). AG News shows smaller gaps when negation is less central to topic identity (H4).

\subsection{Error snapshots}

Per-family classification reports, confusion matrices, and example errors are stored under \texttt{outputs/full/reports/<dataset>/}. Near-synonymous Banking77 intents and sentiment snippets with ironic or double negation remain common failure modes.
